\documentclass[a4paper,11pt]{article}

\usepackage{revy}
\usepackage[utf8]{inputenc}
\usepackage[T1]{fontenc}
\usepackage[danish]{babel}
\usepackage{amsmath,amssymb,amsfonts,amsthm,textcomp}


\revyname{MatematikRevyen}
\revyyear{2026}
% HUSK AT OPDATERE VERSIONSNUMMER
\version{0.1}
\eta{$3-4$ minutter}
\status{Færdig}

\title{Didaktikdukken}
\author{Carl '21}

\begin{document}
\maketitle

\begin{roles}
\role{S1}[Skuespiller] Studerende
\role{S2}[Skuespiller] Studerende
\role{S3}[Skuespiller] Studerende
\role{G}[Skuespiller] Gym-elev
\end{roles}

\begin{props}
\item Brugsanvisning / Manual
\item Didaktikbog
\item Telefon
\item Powerbank
\end{props}


\begin{sketch}
\scene{S1 og S2 står scenen. S3 kommer ind med G med brugsanvisningen}

\says{S3} Hey, undskyld jeg kommer for sent. Jeg skulle lige hente vores eksamensopgaven fra Britta (peger på G)

\scene{S1 og S2 er forvirrede}

\says{S2}[peger på G] Hvem er det?

\says{S3} Det er vores gym-elev

\says{S1} Okay... Men hvorfor er den her?

\says{S3} Britta har jo ændret eksamensformen i didaktik i år. Hver gruppe får en gymnasieelev, som vi skal didaktik...didak...didakttere (?). Følger I hende ikke på LinkedIn??

\says{S2}[stirrer tomt på S3] .... nej 

\says{S1}[til G] Hvad hedder du??

\says{S1}[prikker til G] Hvordan virker den?

\says{S3}[vifter S1 væk] Stop med det der. Det er en rigtig gym-elev.

\says{S3}[tager manualen frem] Vi har fået en brugsanvisning med. 

\says{S3}[læser op] "Guide til gym-elever. Det er ikke unormalt at jeres gym-elev er lidt genert. For at aktivere, prøv følgende: Nævn Mr. Beast"

\scene{G kigger træt på S3, som om det er et dumt spørgsmål}

\says{S3}[læser videre] "Hvis det ikke virker.. bla bla... yndlings tiktok?.. hmm.... Ah! Hvis alt andet fejler så kan man lave et hard reset ved at sige tallene 6, 7 på engelsk."

\says{S1} Altså... six-seven?

\says{G}[laver bevægelsen] Åååhhhh, six-seeeeven

\says{S2} Den virker!

\says{S1}[tager manualen] Hvad skal vi nu?

\says{S1}[læser op] "Afgør niveauet ved at stille basale matematik spørgsmål"

\says{S2}[til G] Hvad er arealet af en cirkel?

\says{G} $\pi r^2$

\says{S2}[overvejer og nikker] Hvad er rumfanget af en kugle?

\says{G} $\tfrac{4}{3}\pi r^3$

\says{S2}[kigger på de andre] ...er det rigtigt?

\scene{S1 og S3 ved det heller ikke}

\says{S1}[læser videre] Anyway,  "Hvis jeres elev er dygtig kan I prøve at lære den nogle logiske udsagn"

\says{S2} Den virker dygtig

\says{S3}[griner drillende af S2] Ja, dygtigere end dig haha

\says{S2} Bro (m/k), du kunne da heller ikk..

\says{S1}[afbryder] Venner, ikke foran G

\says{S1} Okay, G. Hvis $p$ er sandt og $q$ er falskt, hvad er udtrykket $p \Rightarrow q$ så?

\says{G}[ser skræmt ud] Øhhh?

\says{S3} Måske vi skulle prøve med nogle didaktiske metoder

\scene{S3 finder en didaktikbog frem og slår op i den}

\says{S3} "Gør det relaterbart"

\says{S2} Den har jeg!

\says{S2}[til G] Okay, forestil dig at du har doomscrollet i en times tid

\says{G}[Nikker ivrigt genkendende] Mmmh, jaer

\says{S2} Du er træt, du burde sove... Men så! hvis $p$ er sandt og $q$ er falskt, hvad er udtrykket $p \Rightarrow q$ så?

\says{S1}[slår S2 i baghovedet] Ikke sådan din bums. 

\says{S1}[til G] Hvis du har morgenmodul og du sover til kl 10. Er det så sandt at morgenmoduler får dig til at stå tidligt op?

\says{G} Nej, det har sku aldrig været sandt

\says{S3}[begejstret] Ja!!

\says{S3} og så hvis morgenmodul er $p$ og stå tidligt op er $q$, medfører $p$ så $q$?

\says{G}[tøvende] Nej?

\scene{S1-3 er hyped}

\says{S2} Daaaamn!!! Der blev du didaktikeret!

\scene{G forstår ikke helt hvad der sker og bliver lidt genert}

\says{S1}[falder lidt til ro] Okay, okay

\says{S3} Hvad skal vi mere lære den?

\scene{S1 og S3 går lidt til siden og kigger i manualen. G er stadig lidt genert. G tager sin telefon frem og begynder at scrolle. S2 prøver uden held at få kontakt til G}

\says{S2} Øhh venner? Hvad gør man, hvis den begynder at scrolle?

\says{S1}[bladrer i manualen] Ahh her: "Typiske fejl: Din elev er begyndt at scrolle"

\says{S1} "Det kan være meget svært at stoppe en god scroll. Afgør først elevens rps"

\scene{G sætter farten op på sin scroll}

\says{S1}[forvirret] Elevens rps?

\says{S2} Reals per second?

\says{S3} Ahh, ja

\says{S3}[tjekker rps] Hmm, den er vist omkring 2

\scene{S2 prøver stadig at få kontakt til G. G bruger nu flere forskellige fingre til at scrolle med i flydende zen-agtige bevælgeser}

\says{S1}[læser videre] "Er elevens rps over 1.5, har eleven opnået flow-state, og det er meget farligt at stoppe. Vent på den løber tør for strøm. Fratag eventuelle powerbank den måtte have på sig."

\scene{S2 fjerner en powerbank fra G}

\says{S3} Men altså, kan vi ikke bare tage telefonen?

\scene{S3 går hen og prøver at tage telefonen ud af hånden på G}

\says{S1} Nej, vent!

\scene{S3 lykkes med besvær at få telefonen. Telefonen er som en varm bolle i hænderne i S3}

\says{S3} Av! Den er varm

\scene{G står lidt og svajer og falder så død om. S1-3 kigger overraskede og nervøse på hinanden og så ud på publikum. }

\scene{Lys ned}




\end{sketch}
\end{document}
